% Ch2 WS2 -- nomenclature drill. Block 3. Volume matters here: naming is
% a fluency skill, so this sheet is deliberately long and repetitive.
\documentclass{apchem-worksheet}

\apcunitword{Chapter}
\apcunit{2}
\apcunittitle{Atoms, Molecules, and Ions}
\apcdoctitle{Worksheet 2 \textbullet\ Nomenclature Drill}
\apcsubtitle{Zumdahl \S2.8 \textbullet\ fluency sheet --- do it twice, timed the second time}

\begin{document}
\wsheader[Write the compound TYPE beside each answer (I ionic-fixed,
II ionic-variable, III covalent, A acid). Naming errors are almost always
type-identification errors.]

\problem[]{Name each ionic compound (Type I --- fixed-charge metal):}
\begin{parts}
  \item \ce{KBr}: \answerblank[5.5cm]{potassium bromide}
  \item \ce{MgS}: \answerblank[5.5cm]{magnesium sulfide}
  \item \ce{Al2O3}: \answerblank[5.5cm]{aluminum oxide}
  \item \ce{Ca3N2}: \answerblank[5.5cm]{calcium nitride}
  \item \ce{ZnCl2}: \answerblank[6.9cm]{zinc chloride (no numeral needed)}
\end{parts}

\problem[]{Name each (Type II --- variable-charge metal). Deduce the charge
from the anions first.}
\begin{parts}
  \item \ce{FeCl2}: \answerblank[5.5cm]{iron(II) chloride}
  \item \ce{Cu2O}: \answerblank[5.5cm]{copper(I) oxide}
  \item \ce{SnO2}: \answerblank[5.5cm]{tin(IV) oxide}
  \item \ce{PbS}: \answerblank[5.5cm]{lead(II) sulfide}
  \item \ce{Cr2S3}: \answerblank[5.5cm]{chromium(III) sulfide}
\end{parts}

\problem[]{Name each compound containing a polyatomic ion:}
\begin{parts}
  \item \ce{Na2SO4}: \answerblank[5.5cm]{sodium sulfate}
  \item \ce{Mg(OH)2}: \answerblank[5.5cm]{magnesium hydroxide}
  \item \ce{Fe(NO3)3}: \answerblank[5.5cm]{iron(III) nitrate}
  \item \ce{(NH4)2CO3}: \answerblank[5.5cm]{ammonium carbonate}
  \item \ce{KH2PO4}: \answerblank[6.7cm]{potassium dihydrogen phosphate}
  \item \ce{Ca(ClO)2}: \answerblank[5.5cm]{calcium hypochlorite}
\end{parts}

\problem[]{Name each binary covalent compound (Type III):}
\begin{parts}
  \item \ce{SO3}: \answerblank[5.5cm]{sulfur trioxide}
  \item \ce{N2O4}: \answerblank[5.5cm]{dinitrogen tetroxide}
  \item \ce{PCl5}: \answerblank[5.5cm]{phosphorus pentachloride}
  \item \ce{CCl4}: \answerblank[5.5cm]{carbon tetrachloride}
\end{parts}

\problem[]{Name each acid:}
\begin{parts}
  \item \ce{HBr(aq)}: \answerblank[5.5cm]{hydrobromic acid}
  \item \ce{HNO3(aq)}: \answerblank[5.5cm]{nitric acid}
  \item \ce{H2SO3(aq)}: \answerblank[5.5cm]{sulfurous acid}
  \item \ce{H3PO4(aq)}: \answerblank[5.5cm]{phosphoric acid}
\end{parts}

\problem[]{Write the formula for each name:}
\begin{parts}
  \item barium hydroxide: \answerblank[3.5cm]{\ce{Ba(OH)2}}
  \item manganese(II) nitrate: \answerblank[3.5cm]{\ce{Mn(NO3)2}}
  \item diphosphorus pentoxide: \answerblank[3.5cm]{\ce{P2O5}}
  \item ammonium phosphate: \answerblank[3.5cm]{\ce{(NH4)3PO4}}
  \item hydrofluoric acid: \answerblank[3.5cm]{\ce{HF(aq)}}
  \item sodium hypochlorite: \answerblank[3.5cm]{\ce{NaOCl}}
  \item copper(II) sulfate: \answerblank[3.5cm]{\ce{CuSO4}}
  \item potassium bromate: \answerblank[3.5cm]{\ce{KBrO3}}
\end{parts}

\problem[]{Error hunt --- each name below is wrong. Give the correct one and
say which rule was broken.}
\begin{parts}
  \item \ce{FeO} called ``iron oxide''
        \answerlines[2]{Iron(II) oxide --- Fe is variable-charge, so the
        Roman numeral is required.}
  \item \ce{Mg3N2} called ``trimagnesium dinitride''
        \answerlines[2]{Magnesium nitride --- prefixes are for covalent
        compounds only; ionic formulas come from charge balance.}
  \item \ce{Cu(NO3)2} called ``copper(II) nitride''
        \answerlines[2]{Copper(II) nitrate --- \ce{NO3-} is nitrate;
        nitride is \ce{N^3-}.}
\end{parts}

\begin{spiralstrip}{Chapter 2 \textbullet\ blocks 1--2}
  \item Neutrons in \ce{^{108}Ag}: \answerblank[2cm]{61}
  \item Which law does \ce{H2O} always being 88.8\% O illustrate?
        \answerblank[4cm]{definite proportion}
  \item Electrons in \ce{Ti^4+}: \answerblank[2cm]{18}
  \item Charge predicted for a group 6A ion: \answerblank[2cm]{$2-$}
  \item Experiment that showed the atom is mostly empty space:
        \answerblank[3cm]{gold foil}
\end{spiralstrip}

\end{document}
